\chapter{Embedding and Positional Encoding}
\label{ch:embedding}

In this chapter, we implement\keyterm{Embedding (embedding)}, which converts token IDs into dense vectors, and\keyterm{Positional encoding (positional encoding)}, which injects sequence order information into the model.

\section{Why do you need embedding?}

In Chapter 3, the tokenizer changed ``Hello'' to the integer ID 247. However, if you put this integer directly into a neural network, a problem arises. First, there is no guarantee that the integers 247 and 248 are ``similar words.'' ``cat'' might be 247, ``dog'' might be 248, and ``zebra'' might be 248. Second, operations that multiply or add integers have no meaning. ``cat + 1 = dog'' does not hold.

The solution is to convert each token ID to\keyterm{dense vector}. When the vocabulary size is$V$and the embedding dimension is$d$, the embedding is a learnable matrix of size$V \times d$$E$. When token ID$i$comes in, the$i$th row of$E$is retrieved. This row vector becomes the ``semantic representation'' of the token.

\begin{infobox}[Why ``dense''?]
One-hot encoding represents tokens as sparse vectors of vocabulary size$V$. For example, for a vocabulary of 10,000, ID 247 is a vector consisting of 246 0s, 1 1, and the remaining 0s. Multiplying this by$E$is the same as getting the 247th row of$E$. Embedding is an efficient implementation that performs this ``one-hot$\times$matrix'' with a single lookup operation.

As learning progresses, words with similar meanings become closer in the vector space. ``cat'' and ``dog'' are close, and ``cat'' and ``car'' are far. This is the magic of embeddings: before training, they are random vectors, but after training, they become meaningful coordinates.
\end{infobox}

\section{Mathematical definition of embedding}

Given an embedding matrix$E \in \mathbb{R}^{V \times d}$, the embedding of token ID$i$is$E[i] \in \mathbb{R}^d$. This is simply an operation that selects the$i$th row of$E$, and is differentiable (when backwards, only the selected row receives the gradient).

For batch input$\mathbf{x} \in \mathbb{Z}^{B \times L}$(batch$B$, sequence length$L$), the embedding output will be$\mathbf{X} \in \mathbb{R}^{B \times L \times d}$.

\section{Token Embedding avatar}

PyTorch's\code{nn.Embedding}performs exactly this lookup operation. We wrap this thinly and add initialization logic.

\begin{lstlisting}[language=Python]
import math
import torch.nn as nn

class TokenEmbedding(nn.Module):
    def __init__(self, vocab_size: int, d_model: int, pad_id: int = 0):
        super().__init__()
        self.vocab_size = vocab_size
        self.d_model = d_model
        self.pad_id = pad_id
        self.embedding = nn.Embedding(vocab_size, d_model, padding_idx=pad_id)
        self._init_weights()

    def _init_weights(self):
        # normal distribution N(0, 1/sqrt(d_model)) reset (GPT-2 styles)
        nn.init.normal_(self.embedding.weight, mean=0.0,
                        std=1.0 / math.sqrt(self.d_model))
        # Padding token to 0
        with torch.no_grad():
            self.embedding.weight[self.pad_id].fill_(0.0)

    def forward(self, token_ids: torch.Tensor) -> torch.Tensor:
        # [batch, seq] -> [batch, seq, d_model]
        return self.embedding(token_ids)
\end{lstlisting}

\begin{notebox}[Importance of initialization]
In deep learning, weight initialization determines the success or failure of learning. If it is too large, the activation value explodes, and if it is too small, the gradient disappears. GPT-2 uses a normal distribution with standard deviation$1/\sqrt{d\_{model}}$, which is a reasonable choice so that the norm of the initial embedding vector is approximately 1.

Why$1/\sqrt{d}$?  When each element of a$d$-dimensional vector is$\mathcal{N}(0, \sigma^2)$, the expected value of the vector norm is$\sqrt{d} \cdot \sigma$.  If$\sigma = 1/\sqrt{d}$, norm is near 1. This way, the scale of the embedding vector is similar to the layer output, making learning stable.
\end{notebox}

\subsection{padding\_idxrole of}

The\code{padding\_idx}parameter of\code{nn.Embedding}always keeps the embedding of the padding token at 0. If you do this:
\begin{itemize}
\item The output of the padding position becomes 0 and does not affect attention calculation.
\item When going backward, the gradient of the padding row becomes 0 and the weight is not updated.
\end{itemize}

\section{Positional Encoding}

Transformer processes sequences using only attention, but attention itself has no concept of ``order''. Even if you mix the inputs, the output (rearranged) is the same. This is called\keyterm{Order invariance (permutation invariance)}. In natural language, order determines meaning (``dog bites man'' vs ``man bites dog''), so order information must be explicitly injected.

\subsection{intuition: Why Order Matters?}

Attention learns ``how much attention each token will pay to other tokens.'' However, there is no distinction between the ``first token'' and the ``last token''. In the sentence ``I ate rice'', ``I ate'' must be connected to the subject ``I'', but it is difficult for attention to learn this connection without the information that ``the 5th token sees the 1st token''.

Positional encoding adds a unique ``position vector'' to each position, allowing the model to distinguish between positions.

\subsection{Sinusoidal Position encoding (Vaswani et al. 2017)}

Method proposed in the original Transformer paper. Encode the position with sine and cosine functions.

\begin{formula}[Sinusoidal Positional Encoding]
\[
PE\_{(pos, 2i)} = \sin\left(\frac{pos}{10000^{2i/d\_{model}}}\right)
\]
\[
PE\_{(pos, 2i+1)} = \cos\left(\frac{pos}{10000^{2i/d\_{model}}}\right)
\]
Where$pos$is the location and$i$is the dimension index. Even dimensions are sin, odd dimensions are cos.
\end{formula}

The advantage of this method is that there are no learnable parameters. Additionally, a linear relationship for relative positions is established, making it easy for the model to learn.  Since$\sin(\alpha + \beta) = \sin\alpha\cos\beta + \cos\alpha\sin\beta$, the encoding at position$pos + k$is expressed as a linear transformation of the encoding at position$pos$.

\begin{lstlisting}[language=Python]
class PositionalEmbedding(nn.Module):
    def __init__(self, d_model: int, max_len: int = 512, mode: str = "learned"):
        super().__init__()
        self.mode = mode
        if mode == "sinusoidal":
            pe = torch.zeros(max_len, d_model)
            position = torch.arange(0, max_len).unsqueeze(1).float()
            div_term = torch.exp(
                torch.arange(0, d_model, 2).float() * (-math.log(10000.0) / d_model)
            )
            pe[:, 0::2] = torch.sin(position * div_term)
            pe[:, 1::2] = torch.cos(position * div_term)
            self.register_buffer("pe", pe.unsqueeze(0))  # [1, max_len, d_model]
        else:  # learned
            self.pe = nn.Embedding(max_len, d_model)
            nn.init.normal_(self.pe.weight, mean=0.0, std=1.0/math.sqrt(d_model))

    def forward(self, seq_len: int) -> torch.Tensor:
        # [1, seq_len, d_model] return
        if seq_len > self.max_len:
            raise ValueError(f"Sequence length {seq_len} exceeds max_len {self.max_len}")
        if self.mode == "sinusoidal":
            return self.pe[:, :seq_len, :]
        else:
            positions = torch.arange(seq_len, device=self.pe.weight.device)
            return self.pe(positions).unsqueeze(0)
\end{lstlisting}

\begin{figure}[htbp]
\centering
\begin{tikzpicture}[
  scale=0.7,
  every node/.style={font=\small\sffamily}
]
  % sin waveform
  \draw[inkblue, line width=0.5pt, ->] (0, 0) -- (8, 0) node[right] {pos};
  \draw[inkblue, line width=0.5pt, ->] (0, -1.2) -- (0, 1.2) node[above] {value};
  \draw[inkred, line width=1pt, domain=0:7, samples=50] plot (\x, {sin(\x*50)});
  \node[inkred, font=\scriptsize] at (7, 0.8) {dim 0 (sin)};
  \draw[inkblue, line width=1pt, domain=0:7, samples=50] plot (\x, {cos(\x*50)});
  \node[inkblue, font=\scriptsize] at (7, -0.8) {dim 1 (cos)};
\end{tikzpicture}
\caption{Sinusoidal First two dimensions of position encoding. sinclass cosIt changes periodically depending on location.}
\label{fig:sinusoidal}
\end{figure}

\subsection{Learnable location embeddings (GPT-2 styles)}

GPT-2, GPT-3, etc. use learnable embedding vectors for each location. This method is more flexible when there is sufficient data, but it is difficult to respond to lengths that have not been seen during training (e.g., 1024 during training, but 2048 during inference).

\begin{notebox}[Sinusoidal vs Learned]
\begin{itemize}
\item\textbf{Sinusoidal}: No parameters, extrapolation possible (corresponds to sequences longer than the length seen during learning). However, actual performance tends to be slightly lower than that of learnable methods.
\item\textbf{Learned}: Slightly better performance, but fixed length. Extrapolation difficult.
\end{itemize}
Newer models combine the advantages of both methods using Rotary Position Embedding (RoPE). This is covered in Volume 2, Chapter 5.
\end{notebox}

\section{Embedding composition}

Token embeddings and position embeddings are added to create the final input representation. Vaswani et al. suggested scaling the embedding by multiplying it by$\sqrt{d\_{model}}$to keep the scale similar to the positional embedding.

\begin{lstlisting}[language=Python]
class EmbeddingStack(nn.Module):
    def __init__(self, vocab_size, d_model, max_len=512,
                 pad_id=0, pos_mode="learned", dropout=0.1):
        super().__init__()
        self.token_emb = TokenEmbedding(vocab_size, d_model, pad_id=pad_id)
        self.pos_emb = PositionalEmbedding(d_model, max_len, mode=pos_mode)
        self.dropout = nn.Dropout(dropout)
        self.scale = math.sqrt(d_model)

    def forward(self, token_ids: torch.Tensor) -> torch.Tensor:
        batch, seq = token_ids.shape
        tok = self.token_emb(token_ids) * self.scale  # scaling
        pos = self.pos_emb(seq)
        return self.dropout(tok + pos)
\end{lstlisting}

\begin{figure}[htbp]
\centering
\begin{tikzpicture}[
  every node/.style={font=\small\sffamily},
  box/.style={draw=inkblue, fill=inkblue!8, rounded corners=2pt, minimum width=2.5cm, minimum height=0.7cm},
  sum/.style={draw=inkred, fill=inkred!10, circle, minimum size=0.6cm}
]
  \node[box] (ids) at (0, 0) {token IDs};
  \node[box, below=0.6cm of ids] (tok) {Token Embedding};
  \node[box, right=2cm of ids] (pos) {Position};
  \node[box, below=0.6cm of pos] (posemb) {Positional Embedding};
  \node[sum, below=0.8cm of $(tok)!0.5!(posemb)$] (add) {$+$};
  \node[box, below=0.6cm of add] (out) {Dropout $\to$ output};
  \draw[-{Stealth[length=4pt]}, inkblue] (ids) -- (tok);
  \draw[-{Stealth[length=4pt]}, inkblue] (pos) -- (posemb);
  \draw[-{Stealth[length=4pt]}, inkblue] (tok) -- (add);
  \draw[-{Stealth[length=4pt]}, inkblue] (posemb) -- (add);
  \draw[-{Stealth[length=4pt]}, inkblue] (add) -- (out);
\end{tikzpicture}
\caption{EmbeddingStack structure. Add token embedding and location embedding and then apply dropout..}
\label{fig:embedding-stack}
\end{figure}

\subsection{Why do you need scaling??}

The token embedding is initialized to$\mathcal{N}(0, 1/d)$, so the variance of the elements is$1/d$. The elements of the positional embedding (sinusoidal) are in the range$[-1, 1]$and have a variance of approximately$1/2$. When added without scaling, the location information overwhelms the token information.

Multiplying by$\sqrt{d}$makes the variance of the token embeddings$1$, which is balanced with the positional embeddings.

\section{Layer Normalization introduction}

\keyterm{Layer normalization (Layer Normalization)}is used to stabilize learning in Transformer. It will also appear in Attention in the next chapter, so it is introduced here.

\begin{formula}[Layer Normalization]
\[
\text{LN}(x) = \gamma \cdot \frac{x - \mu}{\sqrt{\sigma^2 + \epsilon}} + \beta
\]
Where$\mu, \sigma^2$is the mean and variance calculated along the last dimension.$\gamma, \beta$is a learnable parameter.
\end{formula}

In PyTorch,\code{nn.LayerNorm(d\_model)}is used as a single line. Volume 2 also covers RMSNorm, a variant of LayerNorm.

\begin{infobox}[LayerNorm vs BatchNorm]
BatchNorm normalizes along the batch dimension, so there is a problem with different statistics during learning and inference. LayerNorm operates regardless of batch size by normalizing within each sample. This is why LayerNorm is used in LLM: it works stably even in inference with a batch size of 1.
\end{infobox}

\section{test}

The core tests of the embedding module are checking the output shape and gradient flow.

\begin{lstlisting}[language=Python]
def test_token_embedding_shape():
    emb = TokenEmbedding(vocab_size=100, d_model=32)
    ids = torch.randint(0, 100, (2, 8))  # [batch=2, seq=8]
    out = emb(ids)
    assert out.shape == (2, 8, 32)

def test_positional_sinusoidal():
    pe = PositionalEmbedding(d_model=32, max_len=64, mode="sinusoidal")
    out = pe(seq_len=16)
    assert out.shape == (1, 16, 32)

def test_embedding_stack():
    stack = EmbeddingStack(vocab_size=100, d_model=32, max_len=64, dropout=0.0)
    ids = torch.randint(0, 100, (2, 8))
    out = stack(ids)
    assert out.shape == (2, 8, 32)

def test_positional_exceeds_max_len_raises():
    pe = PositionalEmbedding(d_model=32, max_len=32, mode="learned")
    with pytest.raises(ValueError):
        pe(seq_len=64)  # max_len over
\end{lstlisting}

\section{practice: Embedding visualization}

If you project the embeddings learned from\code{scripts/tiny\_train.py}into 2D using PCA, you can see how the words are distributed. After sufficient learning, you can see that similar words are grouped close together.

\section{summation}

\begin{itemize}
\item Embedding is a learnable lookup table that converts token IDs into dense vectors.
\item initialization uses$N(0, 1/\sqrt{d\_{model}})$, and the padding token is set to 0.
\item Positional encoding injects ordering information into the transformer. There are two methods: sinusoidal and learned.
\item Add token embeddings and position embeddings and apply dropout to create the final input representation.
\item $\sqrt{d\_{model}}$Balance the variance of the two embeddings by scaling.
\item LayerNorm is a batch size-independent regularization that is essential for stabilizing LLM learning.
\end{itemize}

\section{practice problems}

\begin{enumerate}
\item What happens if the embedding dimension$d$is too small (e.g. 8)? What happens if it's too big (e.g. 4096)?
\item What happens if you change the constant\code{10000}to 100 in sinusoidal positional encoding?
\item How would learning change if we did not multiply the token embeddings by$\sqrt{d\_{model}}$?
\item What happens if\code{max\_len}of the learnable position embedding is set to 1024, learned, and then input a 2048 length sequence during inference?
\item How does the output change if you change the$\epsilon$value of LayerNorm from$10^{-5}$to$10^{-1}$?
\end{enumerate}

In the next chapter, we implement an attention mechanism that takes this embedding as input and learns ``how much attention each token will pay to other tokens.''
